\uuid{0ZJI}
\titre{Changement de variables}
\auteur{hanani ; mzari}
\organisation{exo7}
\licence{CC BY-NC-SA 4.0}
\datecreate{2020-01-20}
\niveau{L1}
\module{Analyse}
\chapitre{Calcul d'intégrales}
\sousChapitre{Primitives diverses}
\difficulte{2}
\theme{Primitives, changement de variables}
\exo7qcmid{lille-2.primitives.024}

\contenu{
\begin{question}
\qclassification{127.11}
\qtags{motcle=changement de variables}
Le changement de variable $\displaystyle u=\frac{1}{x}$ donne :
\begin{answers}
\good{$\displaystyle \int \frac{\sin (1/x)}{x^2}\mathrm{d}x=-\int \sin (u)\mathrm{d}u$.}
\good{$\displaystyle \int \frac{\sin (1/x)}{x^2}\mathrm{d}x=\cos\left(1/x\right)+k$, $k\in \Rr$.}
\bad{$\displaystyle \int \sin (1/x)\mathrm{d}x=\int \sin (u)\mathrm{d}u$.}
\bad{$\displaystyle \int \sin (1/x)\mathrm{d}x=-\cos (u)+k$, $k\in \Rr$.}
\end{answers}
\vskip2mm
\begin{explanations}
Avec $\displaystyle u=\frac{1}{x}$, on aura : $\displaystyle \mathrm{d}u=\frac{-\mathrm{d}x}{x^2}$. D'où
$$\int \frac{\sin (1/x)}{x^2}\mathrm{d}x=-\int \sin (u)\mathrm{d}u=\cos (u)+k=\cos\left(1/x\right)+k,\; k\in \Rr.$$
\end{explanations}
\end{question}
}
